\section{The articulation four-cell}

\begin{proposition}[Articulation cell]\label{thm:articulation-cell}\lean{Obsmetrologia.Articulating.Cell}\leanok\uses{def:utterance, def:articulating}
An \emph{articulation cell} is articulating's codomain over an utterance:
the utterance together with two Boolean status bits,
\texttt{tokenArticulated} and \texttt{claimArticulated}. It introduces no
new semantic datum; it records the two-sided status of the utterance that
articulating has brought into view.
\end{proposition}

\begin{proposition}[The four corners]\label{thm:corner}\lean{Obsmetrologia.Articulating.Corner}\leanok\uses{thm:articulation-cell}
The \emph{corners} are the four names for the four status pairs in
$\mathrm{Bool}\times\mathrm{Bool}$: \texttt{potential}, \texttt{active},
\texttt{absurd}, and \texttt{stable}. The partition is the two-slot shape of
an utterance read as an incidence object, not an added taxonomy.
\end{proposition}

\begin{proposition}[Status correspondence]\label{thm:status-correspondence}\lean{Obsmetrologia.Articulating.Corner.ofStatus}\leanok\uses{thm:corner}
The map \texttt{ofStatus} classifies a token/claim status pair as a corner,
and \texttt{toStatus} reads the named corner back as that pair. This is the
coordinate chart through which the four-corner cell is read.
\end{proposition}

\begin{lemma}[Corners biject with status pairs]\label{thm:corner-iso}\lean{Obsmetrologia.Articulating.Corner.ofStatus_toStatus}\leanok\uses{thm:corner, thm:status-correspondence}
Reading a corner to its status pair and classifying back recovers the same
corner. Together with the converse status round-trip, this says the named
corners are exactly the points of $\mathrm{Bool}\times\mathrm{Bool}$.
\end{lemma}

\begin{theorem}[There are exactly four cells]\label{thm:corner-card}\lean{Obsmetrologia.Articulating.Corner.card_eq_four}\leanok\uses{thm:corner, thm:corner-iso}
The list of all corners has length four. The cardinality is forced by the
two independent articulation statuses: two sides give four corners.
\end{theorem}

\begin{proposition}[The corner of a cell]\label{thm:articulation-cell-corner}\lean{Obsmetrologia.Articulating.Cell.corner}\leanok\uses{thm:articulation-cell, thm:corner, thm:status-correspondence}
Every articulation cell has a unique corner, obtained by classifying its two
status bits. This is the projection from articulating's codomain to the
four-corner incidence base.
\end{proposition}

\begin{proposition}[The potential presentation]\label{thm:articulation-cell-initial}\lean{Obsmetrologia.Articulating.Cell.initial}\leanok\uses{thm:articulation-cell-corner}
For an utterance $u$, the initial cell $\texttt{initial}\,u$ has both status
bits false. It is the potential corner over $u$, the presentation from which
predicating is allowed to depart.
\end{proposition}

\begin{proposition}[No empty cell]\label{thm:corner-surjective}\lean{Obsmetrologia.Articulating.Cell.corner_surjective}\leanok\uses{thm:articulation-cell-corner, thm:corner-iso}
Over any utterance, every corner is realized by some articulation cell. The
four-corner partition is therefore inhabited at every point of the utterance
base.
\end{proposition}
